\section{E4: a conditional ranking-tail screen}
\label{sec:e4}
\subsection{Candidate explanation and established mathematics}
The bilinear term's small average energy in E1 does not logically guarantee
accurate Gaussian tail probabilities. Ranking changes can depend on small
margins, and rare-event probabilities may differ even when a variance
component is modest. E4 tests this narrower possibility rather than assuming
the energy result has settled all consequences of non-Gaussian noise.

For a fixed user vector $p$, item pair $i,j$, and signal factors, let
$v=A_i-A_j$, $u=Bp$, $s_x=\sqrt{2}\tau$, and $s_y=\tau\norm{p}$.
The noisy margin has the conditional distribution
\begin{equation}
 L=(v+s_xX)^T(u+s_yY),\qquad X,Y\stackrel{\mathrm{ind}}{\sim}\mathcal N(0,I_r).
 \label{eq:tailmargin}
\end{equation}
Set $a=v/s_x$, $b=u/s_y$, and $k=s_x^2s_y^2$. With independent
noncentral chi-square variables $U,V$ having $r$ degrees of freedom and
noncentralities $\norm{a+b}^2/2$, $\norm{a-b}^2/2$,
\begin{equation}
 L\ \overset{d}{=}\ \frac{\sqrt{k}}{2}(U-V).
 \label{eq:chisquare}
\end{equation}
This is an application of known distributional results
\citep{gaunt2024product}. For a positive clean margin $m=v^Tu$, numerical
integration evaluates $\Pr(U\leq V)$. It is compared with
$\Phi(-m/\sqrt{V_{\mathrm{lin}}+rk})$, where
$V_{\mathrm{lin}}=s_x^2\norm{u}^2+s_y^2\norm{v}^2$, and with the
linearized approximation $\Phi(-m/\sqrt{V_{\mathrm{lin}}})$.

\subsection{Synthetic checks and real-state selection}
Six synthetic cases use 200,000 direct draws each to check moments and
numerical integration, along with rank-one closed forms, symmetry, and
rotation/positive-user-scale invariance. A declared rank-one example with
$a=b=4$ has bilinear variance share $1/33=3.03$\%. Its sign-flip probability
is approximately $0.0000633405$, versus $0.0026743853$ for the variance-
matched Gaussian, a factor of about 42.2. This construction shows why a
tail check is worthwhile; it is not fitted evidence about the trained models.

The real-state screen uses 40 E1 ML-100K Two checkpoints (four ranks, two
privacy labels, five seeds) and ten E3 ML-1M Two-r8 checkpoints. Sixteen
preselected users per dataset and four training-unseen score-rank pairs,
$(10,11),(1,11),(10,100),(1,100)$, yield 3,200 records. The first two pair
types are primary, giving 1,600 primary records. No model is trained and
no held-out relevance label is read. The declared gate requires matched-
Gaussian absolute error above $0.01$ on at least 5\% of primary records
in both datasets.

\subsection{Outcome and numerical scope}
The maximum absolute matched-Gaussian error over all 3,200 records is
$7.66859\times10^{-6}$; the maximum linearized-Gaussian error is
$7.26162\times10^{-5}$. Neither dataset has a primary record above the
$0.01$ threshold. The largest matched error is less than $0.000767$
percentage points. The proposed advancement gate therefore fails.
Group-mean pairwise bilinear variance shares range from approximately
$0.0011$\% to $0.0750$\%, much smaller than the constructed example.
These shares concern a particular user's item-pair margin and are not
the whole-matrix energy fractions reported in E1.

\reportfigure{09_ranking_tail_screen.pdf}{The rejected ranking-tail candidate}
{Maximum absolute flip-probability error for variance-matched and linearized
Gaussian approximations against numerical evaluation of the known conditional
law. Each bar covers 320 state-user-pair records. The logarithmic axis marks
the 0.01 per-pair gate. Even maxima over all four pair types are far below
the threshold. This is a one-step, label-free stability diagnostic, not a
measurement of cumulative recommendation accuracy.}{fig:tails}

There are no tied, zero-user, or numerical-failure records. The largest
reported quadrature error estimate is $2.64884\times10^{-9}$ and the omitted
tail mass is at most approximately $2\times10^{-12}$. A square-root
integration variable avoids the rank-one density singularity. These are
numerical diagnostics, not a rigorous certified error bound for every
possible parameter state. No expanded amplitudes, checkpoints, users,
or pair types were added after inspecting the screen.

The result rejects practically consequential non-Gaussian approximation
error for the sampled one-step round-50 states. It does not prove that all
earlier iterates, all ranks, or multi-round training have Gaussian behavior.
Nor does it refute noisy-score stability results with explicit distributional
assumptions \citep{urmian2026slate}. The correct decision is to retain the
negative result and decline to build a new method around an effect that
was negligible in the tested regime.
\FloatBarrier
